\chapter{Obiettivo e domande di ricerca}
\label{cap:obiettivi}

\section{Obiettivo}

Misurare, con esperimenti ripetibili, come le scelte architetturali (repliche, coda asincrona, orchestrazione) cambiano l'impatto dei guasti, le prestazioni e il recupero di un servizio, e confrontare le misure con previsioni scritte \emph{prima} di misurare.

\section{Domande di ricerca}

\begin{table}[h]
\centering
\caption{Domande di ricerca.}
\label{tab:rq}
\begin{tabularx}{\linewidth}{@{}lL@{}}
\toprule
\textbf{ID} & \textbf{Domanda} \\
\midrule
RQ1 & Quanto downtime e quante richieste fallite produce ogni guasto passando da 1 a N repliche? Quanto si discostano le misure dalle previsioni? Che disponibilità ne risulta con un MTBF ipotizzato? \\
RQ2 & Sincrono contro asincrono, a parità di carico e timeout: come cambiano tempo end-to-end, disponibilità di completamento, ordini persi o duplicati e recupero, sotto picchi e guasti? \\
RQ3 & Quanto ci mette il servizio a tornare effettivamente operativo dopo un guasto, e quanto dipende dai controlli di salute e dal tipo di componente (con o senza stato)? \\
RQ4 & Con parametri allineati, i risultati locali si ripetono su AWS? Dove cambiano e perché? \\
RQ5 & \textit{(opzionale)} Un autoscaling basato sulla lunghezza della coda assorbe meglio un picco? \\
\bottomrule
\end{tabularx}
\end{table}

\section{Fuori ambito}

Portabilità e vendor lock-in; cluster di RabbitMQ e replica del database; multi-zona per i componenti con stato; multi-cloud; ottimizzazione del codice.
